\theabstract{Security observability in Kubernetes-based systems cannot be achieved with traditional methods due to the intensive network communication among microservices and the short-lived nature of workloads. In this thesis, a multi-layer anomaly detection approach based on the joint analysis of network telemetry and runtime events is investigated. Within the scope of the study, network flow data are collected with the eBPF-based Hubble tool and runtime events with the Tetragon tool in a Kubernetes environment; these data are passed through preprocessing and feature extraction steps, and layer-specific unsupervised anomaly detection models are trained. The anomaly scores obtained from the network and runtime layers are combined to produce a unified security score for each pod. The effectiveness of the approach will be examined by validating the single-layer models on publicly available datasets and by evaluating the multi-layer fusion on data generated through controlled attack scenarios.}

\keywords{Anomaly Detection, eBPF, Kubernetes, Machine Learning, Cyber Security.}